\section{Физические механизмы возникновения полярных сияний}

\subsection{Солнечный ветер и межпланетное магнитное поле}

Солнечный ветер --- это поток плазмы из короны Солнца со скоростью 300--800 км/с. На орбите Земли плотность частиц составляет примерно 5--10 см$^{-3}$, температура достигает $10^5$--$10^6$ К. В состав входят электроны, протоны и немного альфа-частиц~\cite{kivelson1995}.

В солнечном ветре «заморожено» межпланетное магнитное поле (IMF) с напряжённостью около 5 нТл. Для взаимодействия с магнитосферой важна южная компонента $B_z$. При $B_z < 0$ создаются условия для пересоединения магнитных линий и усиления сияний.

\subsection{Пересоединение магнитных полей как первичный механизм}

Пересоединение --- процесс разрыва и соединения магнитных линий с изменением их топологии. В магнитосфере Земли он отвечает за передачу энергии от солнечного ветра во внутренние области~\cite{dungey1961}. Энергия, выделяемая при пересоединении, оценивается как:

\begin{equation}
    E_{\text{rec}} = \frac{B^2}{2\mu_0} \cdot V \cdot t,
    \label{eq:energy}
\end{equation}

где $B$ --- напряжённость магнитного поля, $\mu_0$ --- магнитная постоянная, $V$ --- объём области пересоединения, $t$ --- характерное время.

\subsection{Ускорение авроральных электронов}

Существует два основных механизма ускорения электронов до энергий 1--100 кэВ~\cite{lyons1984}.

\subsubsection{Квазистатическое ускорение}

Вдоль магнитных линий могут существовать стационарные электрические поля с напряжением 1--10 кВ. Электроны разгоняются в таких полях, приобретая энергию:

\begin{equation}
    E_e = e \cdot U,
    \label{eq:energy_electron}
\end{equation}

где $e$ --- элементарный заряд, $U$ --- разность потенциалов.

\subsubsection{Альфвеновское ускорение}

Альфвеновские волны --- колебания плазмы вдоль магнитных линий, которые также могут ускорять электроны без постоянного электрического поля. Скорость альфвеновских волн определяется формулой:

\begin{equation}
    v_A = \frac{B}{\sqrt{\mu_0 \rho}},
    \label{eq:alfven}
\end{equation}

где $\rho$ --- плотность плазмы.

Данные спутников THEMIS, Cluster и MMS показывают, что оба механизма работают одновременно~\cite{angelopoulos2008, burch2016}.

\subsection{Возбуждение атмосферных газов и интерпретация спектров}

Ускоренные электроны сталкиваются с атомами кислорода и азота, возбуждая их. Основные эмиссионные линии:
\begin{itemize}
    \item 557,7 нм (зелёная) --- кислород, высоты 100--120 км;
    \item 630,0 нм (красная) --- кислород, высоты 200--300 км;
    \item 391,4 и 427,8 нм (синие/фиолетовые) --- молекулярный азот, ниже 100 км.
\end{itemize}